\subsection{ZERMELO'S THEOREM}

LEMMA 3. Let E be a set, let S be a subset of \(\mathfrak{P}(\mathbf{E})\) , and let \(p\) be a mapping of S into E such that \(p(\mathbf{X}) \notin \mathbf{X}\) for all \(\mathbf{X} \in \mathfrak{S}\) . Then there exists a subset M of E and a well-ordering \(\Gamma\) on M such that, if \(x \leqslant y\) denotes the relation \(y \in \Gamma\langle x \rangle\) and \(\mathrm{S}_x\) denotes the segment {]}←, \(x]\) ,

\begin{enumerate}
\def\labelenumi{(\arabic{enumi})}
\item
  for all \(x \in \mathbf{M}\) we have \(\mathbf{S}_x \in \mathfrak{S}\) and \(p(\mathbf{S}_x) = x\) ;
\item
  \(\mathbf{M} \notin \mathfrak{S}\) .
\end{enumerate}

Let \(\mathfrak{M}\) be the set of subsets \(\mathbf{G}\) of \(\mathbf{E} \times \mathbf{E}\) satisfying the following conditions:

\begin{enumerate}
\def\labelenumi{(\alph{enumi})}
\item
  \(\mathbf{G}\) is the graph of a well-ordering on \(\mathbf{pr}_1\mathbf{G} = \mathbf{U}\) ;
\item
  if \(x \leqslant y\) denotes the relation \((x, y) \in \mathbf{G}\) on \(\mathbf{U}\) , then for each \(x \in \mathbf{U}\) the segment \(S_x\) is such that \(S_x \in \mathfrak{S}\) and \(p(S_x) = x\) . We shall show that if \(G\) and \(G'\) are two elements of \(M\) and if \(U, U'\) denote their first projections, then one of the two sets \(U, U'\) is contained in the other and that if, for example, \(U \subset U'\) , then \(G = G' \cap (U \times U)\) (in other words, that the order relation on \(U\) is induced by the order relation on \(U'\) ) and \(U\) is a segment of \(U'\) .
\end{enumerate}

Consider the set V of elements \(x \in U \cap U'\) such that the segments with endpoint x are the same in U and \(U'\) , and such that the orderings induced on this segment by the orderings on U and \(U'\) are identical. It is clear that V is a segment in both U and \(U'\) and that the orderings induced on V are the same; our assertion will therefore be proved if we show that either V = U or \(V = U'\) . Let us argue by contradiction and suppose that \(V \neq U\) and \(V \neq U'\) . Let x be the least element of U - V in U and let \(x'\) be the least element of \(U' - V\) in \(U'\) ; we have \(V = S_x\) in U, and \(V = S_x'\) in \(U'\) . But by hypothesis, \(V \in S\) and

\[
x = p \left(\mathrm{S} _ {x}\right), \quad x ^ {\prime} = p \left(\mathrm{S} _ {x ^ {\prime}}\right),
\]

so that \(x = x'\) . Hence by definition \(x \in V\) , which is absurd. We may therefore apply Proposition 3 of no. 1 to the set of first projections \(U = pr_1G\) (where \(G \in \mathfrak{M}\) ) and thus obtain a well-ordered set

\[
\mathrm{M} = \bigcup_ {\mathrm{G} \in \mathfrak {M}} \operatorname{pr} _ {1} \mathrm{G}.
\]

It is easily seen that the graph of the ordering on M belongs to M. If we had \(M \in S\) , then, putting \(a = p(M)\) , we should have \(a \notin M\) . We could therefore adjoin to M the element a as greatest element, and the set \(M' = M \cup \{a\}\) would be well-ordered. Since \(M = S_a\) in \(M'\) , we should have \(S_a \in S\) and \(p(S_a) = a\) ; the graph of the ordering on \(M'\) would therefore belong to M, which is absurd.

Note that if \(\phi \notin \mathfrak{S}\) (and in particular if \(\mathfrak{S}\) is empty), the set M whose existence is asserted by Lemma 3 is the empty set; this follows from condition 1 of Lemma 3.

\textbf{THEOREM 1 (Zermelo). Every set E can be well-ordered.}

Let \(\mathfrak{S} = \mathfrak{P}(\mathrm{E}) - \{\mathrm{E}\}\) be the set of all subsets of \(\mathbf{E}\) other than \(\mathbf{E}\) itself. For each \(\mathbf{X} \in \mathfrak{S}\) let \(p(\mathbf{X}) = \tau_x\) ( \(x \in \mathbf{E} - \mathbf{X}\) ); since the relation \(\mathbf{X} \in \mathfrak{S}\) implies \((\exists x)(x \in \mathbf{E} - \mathbf{X})\) , we have \(p(\mathbf{X}) \in \mathbf{E} - \mathbf{X}\) (Chapter I, §4, no. 1) and therefore \(p(\mathbf{X}) \notin \mathbf{X}\) . We may therefore apply Lemma 3, and consequently there exists a well-ordering on a subset \(\mathbf{M}\) of \(\mathbf{E}\) such that \(\mathbf{M} \notin \mathfrak{S}\) ; but the only subset of \(\mathbf{E}\) which does not belong to \(\mathfrak{S}\) is \(\mathbf{E}\) itself, and the theorem is proved.

